\begin{exposition}
Ordinary differential equations are often used to model physical phenomena because many natural systems are governed by relationships between a quantity and its rate of change. In such models, the unknown function represents the state of the system, while the differential equation describes how that state evolves over time. For example, population models describe how a population changes, cooling models describe how temperature changes, motion models describe how position and velocity change, and circuit models describe how current or voltage changes. In each case, the differential equation does not usually give the quantity directly; instead, it gives a rule for how the quantity changes.
\end{exposition}

% ---------------------------------------------------------
% Definition --- Malthusian Population Growth Model
% ---------------------------------------------------------

\begin{definitionbox}{Definition}
\begin{definition}[Malthusian Population Growth Model]
\label{def:ordinary-differential-equations-models-malthusian-population-growth-model}
The \textbf{Malthusian population growth model} assumes that the rate of change of a population is proportional to the current population.

If $P(t)$ denotes the population at time $t$, then the model is
\[
P'(t)=kP(t),
\]
where $k$ is a constant called the \textbf{growth rate} or \textbf{proportionality constant}.
\end{definition}
\end{definitionbox}

\begin{remark*}[Standard quantified statement]
This definition fixes the named kind of differential-equation object by the
conditions stated in the definition.
\end{remark*}

\NoLocalDependencies

\begin{remark*}[Interpretation]
The model says that population growth depends only on the current population size. If $k>0$, the population grows exponentially. If $k<0$, the population decays exponentially. If $k=0$, the population remains constant.
\end{remark*}

\begin{example*}[General Solution]
The general solution of
\[
P'(t)=kP(t)
\]
is
\[
P(t)=Ce^{kt},
\]
where $C$ is a constant. If the initial population is $P(0)=P_0$, then
\[
P(t)=P_0e^{kt}.
\]
\end{example*}






% ---------------------------------------------------------
% Definition --- Decay
% ---------------------------------------------------------

\begin{definitionbox}{Definition}
\begin{definition}[Decay]
\label{def:ordinary-differential-equations-models-decay}
A quantity is said to undergo \textbf{decay} if it decreases over time.

In a differential equation model, if $Q(t)$ denotes the quantity at time $t$, then decay is often modeled by an equation of the form
\[
Q'(t)=-kQ(t),
\]
where $k>0$ is a constant.
\end{definition}
\end{definitionbox}

\begin{remark*}[Standard quantified statement]
This definition fixes the named kind of differential-equation object by the
conditions stated in the definition.
\end{remark*}

\NoLocalDependencies

\begin{remark*}[Interpretation]
The equation
\[
Q'(t)=-kQ(t)
\]
says that the rate of change of the quantity is proportional to the current amount, but with a negative sign. Thus larger amounts decay faster in absolute size, while smaller amounts decay more slowly.
\end{remark*}

\begin{example*}[Exponential Decay]
The general solution of
\[
Q'(t)=-kQ(t)
\]
is
\[
Q(t)=Ce^{-kt}.
\]
If $Q(0)=Q_0$, then
\[
Q(t)=Q_0e^{-kt}.
\]
For $k>0$, this function decreases toward $0$ as $t$ increases.
\end{example*}




% ---------------------------------------------------------
% Definition --- Newton's Law of Cooling
% ---------------------------------------------------------

\begin{definitionbox}{Definition}
\begin{definition}[Newton's Law of Cooling]
\label{def:ordinary-differential-equations-models-newton-s-law-of-cooling}
\textbf{Newton's Law of Cooling} states that the rate at which an object's temperature changes is proportional to the difference between the object's temperature and the surrounding ambient temperature.

If $T(t)$ denotes the temperature of the object at time $t$, and $T_a$ denotes the ambient temperature, then the model is
\[
T'(t)=-k\bigl(T(t)-T_a\bigr),
\]
where $k>0$ is a constant called the \textbf{cooling constant}.
\end{definition}
\end{definitionbox}

\begin{remark*}[Standard quantified statement]
This definition fixes the named kind of differential-equation object by the
conditions stated in the definition.
\end{remark*}

\NoLocalDependencies

\begin{remark*}[Interpretation]
The quantity
\[
T(t)-T_a
\]
measures how far the object's temperature is from the ambient temperature. Newton's Law of Cooling says that the larger this temperature difference is, the faster the object cools or warms toward the ambient temperature.
\end{remark*}

\begin{example*}[General Solution]
The general solution of
\[
T'(t)=-k\bigl(T(t)-T_a\bigr)
\]
is
\[
T(t)=T_a+Ce^{-kt},
\]
where $C$ is a constant. If $T(0)=T_0$, then
\[
T(t)=T_a+(T_0-T_a)e^{-kt}.
\]
\end{example*}
